% ============================================================================
% ARBEITSBLATT: Las asignaturas (Schulfächer)
% ============================================================================
% Sequenz 4: Mi colegio | Block a
% Fachfarbe: #c41e3a (Karminrot)
% Gesamtpunkte: 20
% ============================================================================

\documentclass[11pt, a4paper]{article}

% ============================================================================
% SW-MODUS TOGGLE
% ============================================================================
\newif\ifswmodus
\swmodusfalse

% ============================================================================
% PAKETE
% ============================================================================

\usepackage[utf8]{inputenc}
\usepackage[T1]{fontenc}
\usepackage[ngerman]{babel}
\usepackage{helvet}
\renewcommand{\familydefault}{\sfdefault}

\usepackage[margin=2cm]{geometry}
\usepackage{enumitem}
\usepackage{tabularx}
\usepackage{booktabs}
\usepackage{multirow}

\usepackage{xcolor}
\usepackage{framed}
\usepackage{tcolorbox}
\tcbuselibrary{skins}

\usepackage{fancyhdr}
\usepackage{lastpage}
\usepackage{needspace}

% ============================================================================
% FARBEN
% ============================================================================

\definecolor{spanisch-color}{HTML}{c41e3a}
\definecolor{grau-color}{HTML}{888888}
\definecolor{hellgrau-color}{HTML}{f5f5f5}
\definecolor{orange-color}{HTML}{b35c00}

\definecolor{sw-dark}{HTML}{000000}
\definecolor{sw-gray}{HTML}{666666}
\definecolor{sw-white}{HTML}{ffffff}

\ifswmodus
    \colorlet{spanisch}{sw-dark}
    \colorlet{grau}{sw-gray}
    \colorlet{hellgrau}{sw-white}
    \colorlet{orange}{sw-dark}
\else
    \colorlet{spanisch}{spanisch-color}
    \colorlet{grau}{grau-color}
    \colorlet{hellgrau}{hellgrau-color}
    \colorlet{orange}{orange-color}
\fi

\newcommand{\fachfarbe}{spanisch}

% ============================================================================
% VARIABLEN
% ============================================================================

\newcommand{\thema}{Las asignaturas (Schulfächer)}
\newcommand{\sequenz}{04}
\newcommand{\block}{a}
\newcommand{\fach}{Spanisch}
\newcommand{\klasse}{7}
\newcommand{\maxpunkte}{20}
\newcommand{\datum}{\today}

% ============================================================================
% KOPF- UND FUSSZEILE
% ============================================================================

\pagestyle{fancy}
\fancyhf{}
\renewcommand{\headrulewidth}{0.4pt}
\renewcommand{\footrulewidth}{0.4pt}

\lhead{\fach{} -- Klasse \klasse}
\chead{AB-\sequenz\block}
\rhead{\datum}

\lfoot{}
\cfoot{}
\rfoot{Seite \thepage\ von \pageref{LastPage}}

% ============================================================================
% BEFEHLE
% ============================================================================

\newcommand{\punktebox}[1]{\hfill\fbox{\small \_/#1}}

\newcommand{\luecke}[1]{\rule{#1}{0.4pt}}

\newtcolorbox{merkkasten}[1][]{
    colback=hellgrau,
    colframe=\fachfarbe,
    colbacktitle=white,
    coltitle=\fachfarbe,
    boxrule=1pt,
    arc=0pt,
    left=8pt, right=8pt, top=8pt, bottom=8pt,
    fonttitle=\bfseries,
    attach boxed title to top left={yshift=-2mm, xshift=5mm},
    boxed title style={boxrule=0pt, colback=white},
    title=#1
}

\newtcolorbox{vokabelkasten}{
    colback=white,
    colframe=\fachfarbe,
    boxrule=0.5pt,
    arc=0pt,
    left=5pt, right=5pt, top=5pt, bottom=5pt
}

\newtcolorbox{hinweis}{
    colback=white,
    colframe=orange,
    colbacktitle=white,
    coltitle=orange,
    boxrule=1pt,
    arc=0pt,
    left=5pt, right=5pt, top=5pt, bottom=5pt,
    fonttitle=\bfseries,
    attach boxed title to top left={yshift=-2mm, xshift=5mm},
    boxed title style={boxrule=0pt, colback=white},
    title=Hinweis
}

\newenvironment{aufgabe}[1]{%
    \needspace{5\baselineskip}%
    \nopagebreak[4]%
    \vspace{10pt}
    \noindent\textbf{\textcolor{\fachfarbe}{Aufgabe #1}}
    \nopagebreak[4]%
    \vspace{5pt}
    \nopagebreak[4]%
    \begin{list}{}{\leftmargin=0pt}
    \item[]
}{%
    \end{list}
}

\newlist{teilaufgaben}{enumerate}{2}
\setlist[teilaufgaben,1]{label=\alph*), leftmargin=2em}
\setlist[teilaufgaben,2]{label=\roman*), leftmargin=2em}

\newcommand{\es}[1]{\textcolor{\fachfarbe}{\textbf{#1}}}

% ============================================================================
% DOKUMENT
% ============================================================================

\begin{document}

% ============================================================================
% KOPFBEREICH
% ============================================================================

\begin{center}
    {\Large\bfseries\textcolor{\fachfarbe}{Arbeitsblatt: \thema}}
\end{center}

\vspace{5pt}

\noindent
\begin{tabularx}{\textwidth}{|X|X|X|}
\hline
\textbf{Name:} \luecke{4cm} &
\textbf{Klasse:} \klasse &
\textbf{Datum:} \luecke{2cm} \\
\hline
\end{tabularx}

\vspace{15pt}

% ============================================================================
% MERKKASTEN: Schulfächer
% ============================================================================

\begin{merkkasten}[Merkkasten: Las asignaturas (Schulfächer)]
\textbf{Schulfächer auf Spanisch:}

\vspace{0.5em}

\begin{tabular}{p{5cm}p{5cm}}
\es{las matematicas} = \luecke{3cm} & \es{la historia} = \luecke{3cm} \\[0.8em]
\es{el espanol} = \luecke{3cm} & \es{la geografia} = \luecke{3cm} \\[0.8em]
\es{el ingles} = \luecke{3cm} & \es{la biologia} = \luecke{3cm} \\[0.8em]
\es{el aleman} = \luecke{3cm} & \es{la musica} = \luecke{3cm} \\[0.8em]
\es{el arte} = \luecke{3cm} & \es{la educacion fisica} = \luecke{3cm} \\[0.8em]
\es{la informatica} = \luecke{3cm} & \\
\end{tabular}

\vspace{0.8em}

\textbf{Meinung ausdrücken:}

\vspace{0.3em}

\begin{tabular}{ll}
\es{Me gusta} + \textit{el/la + Fach} & = Das Fach gefällt mir. \\[0.4em]
\es{Me gustan} + \textit{las matematicas} & = Die Fächer gefallen mir. (Plural!) \\[0.4em]
\es{No me gusta(n)} + ... & = ... gefällt/gefallen mir nicht. \\
\end{tabular}

\vspace{0.5em}

\textit{Tipp: Übertrage die Übersetzungen aus dem Lernpfad!}
\end{merkkasten}

\vspace{10pt}

% ============================================================================
% AUFGABE 1: Zuordnung (4 Punkte)
% ============================================================================

\begin{aufgabe}{1}
Verbinde die spanischen Schulfächer mit den deutschen Übersetzungen. \punktebox{4}
\end{aufgabe}

\begin{center}
\begin{tabular}{rcl}
\es{las matematicas} & \luecke{3cm} & Geschichte \\[0.8em]
\es{la historia} & \luecke{3cm} & Sport \\[0.8em]
\es{el ingles} & \luecke{3cm} & Mathematik \\[0.8em]
\es{la educacion fisica} & \luecke{3cm} & Englisch \\
\end{tabular}
\end{center}

\vspace{15pt}

% ============================================================================
% AUFGABE 2: Artikel (4 Punkte)
% ============================================================================

\begin{aufgabe}{2}
Ergänze den richtigen Artikel: \textbf{el} oder \textbf{la}. \punktebox{4}
\end{aufgabe}

\begin{vokabelkasten}
\textit{Tipp:} el = maskulin, la = feminin
\end{vokabelkasten}

\vspace{0.5em}

\noindent
\begin{tabular}{p{7cm}p{7cm}}
a) \luecke{1.5cm} espanol & e) \luecke{1.5cm} biologia \\[0.8em]
b) \luecke{1.5cm} musica & f) \luecke{1.5cm} arte \\[0.8em]
c) \luecke{1.5cm} geografia & g) \luecke{1.5cm} aleman \\[0.8em]
d) \luecke{1.5cm} informatica & h) \luecke{1.5cm} ingles \\
\end{tabular}

\vspace{15pt}

% ============================================================================
% AUFGABE 3: Meinung äußern (6 Punkte)
% ============================================================================

\begin{aufgabe}{3}
Schreibe Sätze über deine Lieblingsfächer und Fächer, die du nicht magst. \punktebox{6}
\end{aufgabe}

\begin{hinweis}
Benutze: \es{Me gusta...} / \es{Me gustan...} / \es{No me gusta...} / \es{No me gustan...}
\end{hinweis}

\vspace{0.5em}

\noindent
\textbf{Beispiel:} \textit{Me gusta el espanol. No me gustan las matematicas.}

\vspace{0.5em}

\noindent
a) Schreibe zwei Fächer, die du magst:

\vspace{0.3em}
\noindent\rule{\textwidth}{0.4pt}\\[0.5em]
\noindent\rule{\textwidth}{0.4pt}

\vspace{0.8em}

\noindent
b) Schreibe zwei Fächer, die du nicht magst:

\vspace{0.3em}
\noindent\rule{\textwidth}{0.4pt}\\[0.5em]
\noindent\rule{\textwidth}{0.4pt}

\vspace{0.8em}

\noindent
c) Frage einen Partner / eine Partnerin und schreibe die Antwort auf:

\vspace{0.3em}
\noindent
\textit{?Te gusta...?} \luecke{10cm}

\vspace{0.3em}
\noindent\rule{\textwidth}{0.4pt}

\vspace{15pt}

% ============================================================================
% AUFGABE 4: Stundenplan lesen (6 Punkte)
% ============================================================================

\begin{aufgabe}{4}
Lies den Stundenplan von Maria und beantworte die Fragen auf Spanisch. \punktebox{6}
\end{aufgabe}

\begin{center}
\textbf{El horario de Maria (lunes)}

\vspace{0.5em}

\begin{tabular}{|c|l|}
\hline
\textbf{Hora} & \textbf{Asignatura} \\
\hline
8:00 - 9:00 & matematicas \\
\hline
9:00 - 10:00 & espanol \\
\hline
10:00 - 11:00 & ingles \\
\hline
11:00 - 11:30 & \textit{recreo} (Pause) \\
\hline
11:30 - 12:30 & historia \\
\hline
12:30 - 13:30 & educacion fisica \\
\hline
\end{tabular}
\end{center}

\vspace{0.8em}

\noindent
a) ?Que asignatura tiene Maria a las 8:00?

\vspace{0.3em}
\noindent
\textit{Maria tiene} \luecke{6cm}

\vspace{0.8em}

\noindent
b) ?Que asignatura tiene Maria a las 12:30?

\vspace{0.3em}
\noindent
\textit{Maria tiene} \luecke{6cm}

\vspace{0.8em}

\noindent
c) ?Cuantas asignaturas tiene Maria el lunes? (Zahl ausschreiben!)

\vspace{0.3em}
\noindent
\textit{Maria tiene} \luecke{4cm} \textit{asignaturas.}

% ============================================================================
% PUNKTEÜBERSICHT
% ============================================================================

\vspace{20pt}
\hrule
\vspace{10pt}

\begin{flushright}
\textbf{Punkte:} \luecke{1cm} / \maxpunkte
\end{flushright}

\end{document}
